\section{Related Work}

Architectures for video understanding have mirrored advances in image recognition.
Early video research used hand-crafted features to encode appearance and motion information~\cite{laptev_ijcv_2005, wang_ijcv_2013}.
The success of AlexNet on ImageNet~\cite{krizhevsky_neurips_2012, deng_cvpr_2009} initially led to the repurposing of 2D image convolutional networks (CNNs) for video as ``two-stream'' networks~\cite{karpathy_cvpr_2014, simonyan_neurips_2014, ng_cvpr_2015}.
These models processed RGB frames and optical flow images independently before fusing them at the end.
Availability of larger video classification datasets such as Kinetics~\cite{kay_arxiv_2017} subsequently facilitated the training of spatio-temporal 3D CNNs~\cite{carreira_cvpr_2017, feichtenhofer_neurips_2016,tran_iccv_2015} which have significantly more parameters and thus require larger training datasets.
As 3D convolutional networks require significantly more computation than their image counterparts, many architectures factorise convolutions across spatial and temporal dimensions and/or use grouped convolutions~\cite{sun_iccv_2015,tran_iccv_2019,tran_cvpr_2018,xie_s3d_eccv_2018, feichtenhofer_cvpr_2020}.
We also leverage factorisation of the spatial and temporal dimensions of videos to increase efficiency, but in the context of transformer-based models.

Concurrently, in natural language processing (NLP), Vaswani~\etal~\cite{vaswani_neurips_2017} achieved state-of-the-art results by replacing convolutions and recurrent networks with the transformer network that consisted only of self-attention, layer normalisation and multilayer perceptron (MLP) operations.
Current state-of-the-art architectures in NLP~\cite{devlin_naacl_2019, raffel_jmlr_2020} remain transformer-based, and have been scaled to web-scale datasets~\cite{brown_gpt3_neurips_2020}.
Many variants of the transformer have also been proposed to reduce the computational cost of self-attention when processing longer sequences~\cite{child_arxiv_2019, choromanski_arxiv_2020,kitaev_iclr_2020, tay2020long, tay2020efficient, wang_arxiv_2020} and to improve parameter efficiency~\cite{lan2019albert, dehghani2018universal}.
Although self-attention has been employed extensively in computer vision, it has, in contrast, been typically incorporated as a layer at the end or in the later stages of the network~\cite{wang_cvpr_2018, carion_eccv_2020,  huang_ccnet_iccv_2019,wang_vistr_arxiv_2020, zhang_dgmn_cvpr_2020} or to augment residual blocks~\cite{hu_squeeze_excite_cvpr_2018, cao_gcnet_cvprw_2019, chen_danet_neurips_2018, srinivas_cvpr_2021} within a ResNet architecture~\cite{he_cvpr_2016}.

Although previous works attempted to replace convolutions in vision architectures~\cite{parmar_icml_2018, ramachandran_neurips_2019, shen_gsa_iclr_2021}, it is only very recently that Dosovitisky~\etal~\cite{dosovitskiy_iclr_2021} showed with their ViT architecture that pure-transformer networks, similar to those employed in NLP, can achieve state-of-the-art results for image classification too.
The authors showed that such models are only effective at large scale, as transformers lack some of inductive biases of convolutional networks (such as translational equivariance), and thus require datasets larger than the common ImageNet ILSRVC dataset~\cite{deng_cvpr_2009} to train.
ViT has inspired a large amount of follow-up work in the community, and we note that there are a number of concurrent approaches on extending it to other tasks in computer vision~\cite{wang_maxdeeplab_arxiv_2020,wang_pvt_arxiv_2021, zhao_point_transformer_arxiv_2020,zheng_setr_arxiv_2020} and improving its data-efficiency~\cite{touvron_arxiv_2020,pan_hvt_arxiv_2021}.
In particular, \cite{bertasius_arxiv_2021,neimark_arxiv_2021} have also proposed transformer-based models for video.




In this paper, we develop pure-transformer architectures for video classification.
We propose several variants of our model, including those that are more efficient by factorising the spatial and temporal dimensions of the input video.
We also show how additional regularisation and pretrained models can be used to combat the fact that video datasets are not as large as their image counterparts that ViT was originally trained on.
Furthermore, we outperform the state-of-the-art across five popular datasets.


